\section{Arbitrary predicates inside small cells}
Let each cell impose an arbitrary Boolean relation among its entries. We may additionally prescribe the exact total weight $W$. Write $T$ for the number of compatible words on that layer, again with $H=\log\binom NW$.
\begin{theorem}\label{th:relations}
Suppose the largest cell has $o(N)$ entries and $T>0$. For every $0\le\lambda\le1$ there are at least $2^{\lambda\log T-o(N)}$ compatible words with $C(y)\le [H-(1-\lambda)\log T+o(N)]/n$. The losses are uniform given an upper bound tending to zero on the largest cell's fraction of $N$.
\end{theorem}
\begin{proof}
The preceding dictionary only requires a short description of every possible next-block probability bucket. We construct such descriptions even when the predicates themselves are expensive to evaluate.

\paragraph{Spreading the cells.} Let $p$ be the maximum cell fraction. In the array of $N/b$ blocks by $b$ positions, independently translate the high address bits in each position-column by a random XOR shift. Each cell's occupancy of a fixed physical block is a sum of independent indicators of mean at most $pb$. Put $t=\lceil10 n/\log(1/p)\rceil=o(n)$. For large $n$ its upper-tail probability beyond $t$ is at most $(epb/t)^t\le2^{-3n}$. A union bound fixes shifts so every block meets every cell in at most $t$ sites. Computing this permutation or its inverse needs only polynomially many gates in $n$.

\paragraph{Recording global weight.} Put $k=\lfloor\sqrt n\rfloor$. Fix the weights of all cells larger than $k$ to some profile retaining at least $T'=T2^{-o(N)}$ feasible words. Indeed, $\sum_{M_i>k}\log(M_i+1)=O(N\log(k+1)/k)$. The remaining global constraint on the smaller cells is their total residual weight. At a prefix the uniform feasible completion now factors into independent large-cell laws and one law on the small cells jointly.

For each large cell touched in the current block, record the conditional probabilities of its at most $2^t$ local patterns by rounding their negative logarithms upwards to multiples of $1/b$, with a sentinel above $2b$. Any block pattern hitting a sentinel has probability at most $2^{-2b}$. Summing over all block patterns loses at most $2^{-b}$. For the others the products of the rounded probabilities have error at most a factor of two.

For a touched small cell record the set of conditionally permitted completions on its at most $k$ remaining sites, and which sites occur in the current block. Altogether these touched small cells have at most $bk$ remaining sites. Let $W'$ be the residual required weight in all small cells. For each $0\le j\le bk$, record the integer part of the logarithm of the number of compatible assignments of weight $W'-j$ on the untouched small cells, or a zero sentinel. Each integer needs $O(n)$ bits. For a candidate block word, sum over compatible patterns on the rest of the touched small cells, weighting each by the recorded approximation for the needed untouched weight. The resulting mass is within a factor of two of the true numerator, and normalizing distorts probabilities by at most another constant factor. The entire local description, including the arrangement of sites, uses $D=\operatorname{poly}(n)(2^t+2^k)$ bits. This counts possible descriptions. A circuit for a satisfying word uses only its short index.

\paragraph{Using the dictionary.} Discard all additional block words whose estimated probability is less than $2^{-b-C}$ for a sufficiently large constant $C$. At least half the mass survives. Group remaining words by weight and dyadic estimated probability; a heaviest bucket has actual mass $\Omega(1/b^2)$ and its atom probabilities differ by only a fixed factor. Thus (\ref{eq:credit}) still holds. The logarithm of the number of candidate buckets is $O(D+\log n)$. Oversampling by $O(D+n)$ in the hypergeometric argument gives a valid common dictionary with an additional $O(\log(D+n))=O(t+k+\log n)=o(n)$ bits per block. The path count and packed circuit lookup of Section~2 give the theorem.
\end{proof}
The size condition is needed for arbitrary predicates. On a dyadic address subcube of some fixed density $\alpha>0$ exclude all restrictions realizable with at most $\alpha N/(3n)$ gates and leave other sites free. A uniform balanced table matches any specific subcube pattern with probability at most $O(\sqrt N)2^{-\alpha N}$, by conditioning independent bits on total weight. The circuit census implies that the excluded fraction is exponentially small, yet every remaining table needs more than $\alpha N/(3n)$ gates. Quotas retain the short conditional descriptions of Section~2 even in this regime.
